% This module uses amsmath, amssymb, amsthm, and the main document's
% definition, lemma, corollary, and proof environments.
% The normalization lemma below is a new, manually proved interface result.
% The separately cited Tamarin theorem is an existing symbolic result.

\section{Reservation origins, persistent records, and recovery}
\label{sec:security-model}

An interrupted cryptographic operation can retain its random tape before it
has produced a key or a ciphertext.  Its later execution and confirmation
therefore have to be related to the original reservation.  We make this
relationship explicit at two levels: the origin of a cryptographic operation,
and the origin of a saved session-state version.  A third identifier describes
the physical image that carries that version.  This distinction lets us
account for a target record copied by another session's write, an unapplied
pending image, and a ciphertext acquired before its wrapping key.

The model is an interface model for an asynchronous protocol machine
$\mathcal P$.  The adversary schedules public messages, process loss, honest
recovery, and the disclosure operations specified below.  Application
deliveries are private to the honest applications.  Capture grants read access
to the stated assets; it does not grant an additional operation for replacing
the authoritative journal.  An implementation with adversarial rollback or
forking requires those transitions in its storage machine.  The protocol,
its cryptographic primitives, and its public outputs are parameters of the
record-normalization lemma.  Fix finite execution, record-size, and query
budgets; the wrapper's bookkeeping is polynomial in their total
representation size.

Table~\ref{tab:record-notation} summarizes the principal objects.

\begin{table}[t]
\centering
\caption{Principal objects in the record model. The three identities separate
an original random draw, an immutable session-state version and its physical
storage representation.}
\label{tab:record-notation}
\small
\begin{tabularx}{\linewidth}{@{}>{\raggedright\arraybackslash}p{.25\linewidth}X@{}}
\toprule
Notation & Meaning\\
\midrule
$o,\rho_o$ & Original invocation and its reserved random tape\\
$v,X_v$ & Immutable session-state version and its payload\\
$h$ & Physical image or pending object carrying one or more versions\\
$\kappa$ & Identity of actual key material, shared by equal keys\\
$\mathsf{res}(o)$ & Event at which the original random tape is drawn\\
$\mathsf{save}(v)$ & First retention of $v$ on the durable acquisition surface\\
$\mathcal C_t,\mathcal L_t$ & Acquired and readable original versions at prefix $t$\\
\bottomrule
\end{tabularx}
\end{table}

\subsection{Three origins and their event order}

\begin{definition}[Reservation origin]
\label{def:reservation-origin}
A reservation origin is a fresh invocation index $o$, together with
\[
  \bigl(\mathsf{kind}_o,\mathsf{owner}_o,s_o,e_o,B_o,
        \mathsf{res}(o)\bigr).
\]
Here $\mathsf{kind}_o$ distinguishes generation, encapsulation, and any other
reserved operation; $s_o,e_o$ identify its session and epoch; $B_o$ is the full
canonical binding, including the admitted policy and operation inputs; and
$\mathsf{res}(o)$ is the event at which its random tape $\rho_o$ is drawn.
Execution or exact replay of $o$ consumes the same $\rho_o$ and the same
binding.  A new draw receives a new origin even if a failed attempt's external
operation identifier is reused.
\end{definition}

The ideal source in this model draws independent tapes of the prescribed
lengths.  Explicit entropy-disclosure queries reveal designated tapes.
Independence is a distributional assumption about the source at sampling.
Later API calls preserve the sampled tapes. The generation and encapsulation budgets
$G,C$ count original draws, including draws in attempts that subsequently
fail.  Exact retries consume no new slot.  Independently sampled operations
retain distinct origins even if some public outputs happen to coincide.

A record version $v$ has an immutable payload $X_v$ and semantic metadata
\[
  \bigl(\mathsf{subject}(v),\mathsf{role}(v),
         \mathsf{epoch}(v),\mathsf{phase}(v)\bigr).
\]
The subject identifies the session whose state is represented. A write to the
enclosing journal preserves this subject, including a write by another session.
A version can
refer to several reservation origins and to roots or traffic state derived
from earlier origins.  Preparing a version allocates its handle.  The event
$\mathsf{Saved}(v)$ occurs when that version first enters the durable
acquisition surface, and its time is denoted $\mathsf{save}(v)$.  In particular,
an intact target in a durable pending-write intent is saved before the
corresponding image is applied.  A storage operation may have made this
transition even when its caller receives an unknown-outcome error.

A physical object $h$ has a separate creation or sealing index.  It may carry
one version, several sessions' versions, an exact earlier image, or a pending
target.  An unchanged version keeps its $v$ and $\mathsf{save}(v)$ when it is
copied, archived, applied from a pending intent, or included in a newly sealed
aggregate.  A new semantic session-state version receives a new $v$ even when
some of its fields are unchanged.  These rules distinguish the original
reservation from each session-state update that retains it.

The following transition rules specify the provenance discipline.

\begin{enumerate}
  \item \emph{Reserve and execute.}  Reserve creates $o$ at the actual draw.
  Execution uses its retained binding and tape.  Unknown-outcome
  reconciliation keeps that origin; a definite new attempt allocates another.
  \item \emph{Prepare and retain.}  Preparing $v$ fixes $X_v$ and its provenance.
  Retaining the first intact image or pending target containing $v$ emits
  $\mathsf{Saved}(v)$.  An attempt that retains no such object has no successful
  acquisition event.
  \item \emph{Copy, apply, and reseal.}  An exact copy preserves the complete
  object and its sealing origin.  Resealing creates a new sealing invocation
  and its actual new wrapper randomness, while preserving the referenced
  record versions.  Applying an existing pending target copies its retained
  bytes.
  \item \emph{Crash and resume.}  Process loss removes volatile references.
  Recovery reconstructs the retained handles and their original provenance.
  It does not turn a retained reservation into a new draw.
  \item \emph{Confirm and delete.}  Confirmation changes protocol state and may
  authorize particular host-side deletions.  Neither action changes a
  reservation's origin or removes an adversary's previously acquired copy.
\end{enumerate}

For an original reservation actually retained in $v$, causality requires
\begin{equation}
  \mathsf{res}(o)\ \leq\ \mathsf{save}(v).
  \label{eq:origin-causality}
\end{equation}
A later physical object's creation time gives an additional timestamp, not a
replacement for either side of~\eqref{eq:origin-causality}.  The interval
between drawing $\rho_o$ and first retaining it remains explicit.  The
symbolic reservation-and-save transition can represent that interval only
through a separately justified source correspondence.

\subsection{Acquisition, delayed readability, and key identity}

Wrapping-key handles resolve to identities of actual key material.  Two
session labels that designate the same key resolve to the same identity
$\kappa$.  Disclosure through either label therefore enables every acquired
wrapper protected by $\kappa$.  A rotation creates a different identity when
it creates different key material.  Signing-key identities are treated in
the same way across sessions and roles.

For each intact physical object $h$, its construction records a finite set
\[
  \mathsf{paths}(h)=\{(v,\gamma)\},
\]
where $\gamma$ is the set of wrapping-key identities on a recoverable
serialization path to $X_v$.  Multiple paths to one version are permitted.
An unprotected record has $\gamma=\varnothing$.  Nesting a wrapper adds its
key identity; aggregation unions its members' paths.  An exact ciphertext
copy keeps its paths.  Resealing decoded records reconstructs the paths under
the new wrapper and retains the same $v$ values.  Thus resealing decoded
plaintext and wrapping an old ciphertext are distinct operations.

Only intact, recoverably embedded record payloads create these paths.  A
public hash or MAC of a private record is a public computation of that record,
not a path to the record itself.  The renderer introduced below preserves
that computation and its output.  This distinction prevents a public head
query from being relabeled as a snapshot merely because its computation used
private bytes.

\begin{definition}[Record acquisition and readability]
\label{def:record-acquisition}
A successful $\mathsf{Capture}(h)$ returns the actual retained bytes of $h$
and adds an immutable copy to the adversary's archive $\mathcal A_t$.
Let $\mathcal K_t$ be the wrapping-key identities disclosed by time $t$.
Define
\begin{align}
  \mathcal C_t
    &=\{v:\exists h\in\mathcal A_t,\ \exists\gamma,\
                    (v,\gamma)\in\mathsf{paths}(h)\}, \label{eq:acquired-records}\\
  \mathcal L_t
    &=\{v:\exists h\in\mathcal A_t,\ \exists\gamma,\
          (v,\gamma)\in\mathsf{paths}(h)
          \ \wedge\ \gamma\subseteq\mathcal K_t\}.
          \label{eq:readable-records}
\end{align}
Each capture records $\mathsf{AcquireOriginal}(v,t)$ for every version in its
support, including versions already acquired through another path.
An unsuccessful request for unavailable state returns its ordinary error and
adds no acquisition.
\end{definition}

Both sets describe the specified disclosure interfaces, using provenance and
key identities.  They do not test whether the adversary has guessed a secret.
The read set can grow without another capture: if $h$ was obtained at
$t_1$ and its final missing wrapping key is disclosed at $t_2>t_1$, the
corresponding versions enter $\mathcal L_{t_2}$.  Host deletion at an
intermediate time leaves $h$ in $\mathcal A_{t_2}$.  Acquisition, by contrast,
already occurred at $t_1$.

An additional key that is explicitly recoverable from an opened record can be
included by a declared extraction rule.  In this extension, iterate
\eqref{eq:readable-records} and the declared key-extraction rules to their least
fixed point after each event.  Each rule identifies its input bytes, its
deterministic extraction algorithm, and the resulting actual key identity.
Only rules with available inputs fire.  This accounts for nested protection
without assuming that inaccessible keys have already been learned.

Signer-file capture is a separately typed disclosure.  For correspondence
with the retained symbolic theorem, it records exposure of the actual signing
key conservatively at acquisition, as does that theorem's signer-file rule.
The readability ledger can additionally record when the file's contents
become decryptable.  In particular, a wrapping-key disclosure cannot be
scoped to traffic state when the declared capture also includes the matching
signer file.

\paragraph{Constructive exposure.}
From public protocol inputs, explicit leaks, and readable record payloads, we
form a labelled replay-dependency graph.  Its edges are specified algorithms:
unseal a reservation with its matching key; regenerate a key from its original
tape; replay encapsulation on the bound offer; decapsulate a published
ciphertext; derive a root; and derive the direction-specific traffic and ACK
material.  Directly stored roots, skipped keys, and plaintext enter through
their own record fields.  For the implemented root schedule, an edge is
\[
 (R_{e-1},S_e,C_e)\longmapsto
 R_e=\mathsf{HKDF}\bigl(R_{e-1},S_e,
       D_{\rm root}\Vert H_{\rm core}(C_e),32\bigr).
\]
Closing these explicit edges over the nodes of the bounded execution yields
a computable set of message keys with replay witnesses.  It supplies positive
exposure certificates; the
confidentiality game below is the separate criterion for messages outside
the declared disclosure conditions.  Confirmation counters belong to the
public progress state and do not reset this graph.

\subsection{An original-record implementation of the copy interface}

We now isolate a property of record management that is independent of KEM
secrecy.  The following two storage machines use the same protocol machine
$\mathcal P$, random-tape source, scheduler, and cryptographic interfaces.

The \emph{literal-copy machine} materializes the immutable payload $X_v$ in
each indicated serialization position.  Its objects include current images,
pending targets, and retained historical copies.  It records the origin tags
alongside these values as private bookkeeping.

The \emph{original-record machine} keeps one immutable table entry
$\mathcal O[v]=X_v$ and represents each serialization position by a reference
to $v$.  Its wrapper maintains the physical-object directory, pending states,
adversary archive, origin tags, and key identities.  It has the following
specified interfaces.

\begin{enumerate}
  \item $\mathsf{RegisterOriginal}(v,X_v)$ is a private write by the honest
  machine.  It fixes one value at a fresh $v$.  The wrapper receives the
  handle and the provenance metadata.
  \item $\mathsf{RenderOriginal}(\theta)$ evaluates the supplied serialization
  recipe $\theta$ by dereferencing $\mathcal O$.  A recipe specifies literal
  public fields, record references, references to already rendered objects,
  and each sealing, hash, MAC, or other admitted byte-processing call.
  A fresh sealing call uses a fresh invocation index and the same prescribed
  coin distribution as the literal-copy machine.  An exact-object read or
  replay returns its cached original bytes.
  \item $\mathsf{AcquireOriginal}(v,t)$ records acquisition of that version.
  This is an event, not an assertion of decryption.  A capture's externally
  returned bytes come from the rendered object cache.
  \item $\mathsf{RevealOriginal}(v)$ returns the original payload $X_v$.
  The wrapper uses this interface for a plaintext record-opening response
  only after a captured path is enabled by the disclosed keys, and caches
  the answer for subsequent openings.
\end{enumerate}

The renderer is an explicit trusted interface in this intermediate
experiment.  It has access to the original payload table and evaluates the
stated byte operations, including operations on undisclosed payloads.
Consequently, it can produce real ciphertexts, commitment digests, and report
MACs.  This definition does not supply an algorithm for producing those bytes
from an unknown KEM-challenge tape.  It gives a precise location for the
subsequent cryptographic replacement of that interface.

The operational contract has three further requirements.
First, all publicly visible payload-dependent bytes and statuses are obtained
through the specified renderer or through $\mathcal P$'s common protocol
computation.  Memory addresses and host timing are not additional outputs of
this interface.
Second, actual byte-operation order, error responses, partial transaction
outcomes, and cryptographic-oracle calls are retained.  Caching an exact
object does not discard a native hash-validation call, observable retry, or
write attempt: such an action remains a renderer or protocol action.
Third, the reference metadata and the key-identity relation are available to
the wrapper without inspecting a selected hidden reservation's coins.
These are checkable interface conditions; none asserts the secrecy of a
record, a root, or a message.

\begin{lemma}[Alias normalization]
\label{lem:alias-normalization}
For the operations and renderer contract above, every adversary
$\mathcal A$ interacting with the literal-copy machine has an online wrapper
$\mathcal B$ interacting with the original-record machine such that their
public views have the same distribution.  Their reservation-origin ledgers,
record-save events, acquisition events, and read sets are identical after
projection onto the original handles.

The wrapper makes at most $|\mathcal L_T|$ distinct
$\mathsf{RevealOriginal}$ calls by terminal time $T$.  Repeated physical
copies, captures, and openings do not multiply this number.  The wrapper
preserves every original reservation and every fresh sealing invocation, so
it does not merge their respective cryptographic resource counts.

The statement also holds for an adversary with quantum private state and
quantum access to common public cryptographic oracles, provided the storage
interface is classical and both experiments retain the same oracle calls.
\end{lemma}

\begin{proof}
Maintain a map from every literal record occurrence to its immutable handle
$v$, and from every native object identifier to a rendered-object entry.
Couple the protocol random tapes, sealing coins, scheduler choices, and public
cryptographic oracles between the experiments.  The invariant is that
dereferencing every original-record recipe produces the corresponding
literal-copy bytes, with the same semantic provenance and physical
disposition.

For a new record, register its payload once and extend the occurrence map.
For a fresh rendering, submit the same recipe and byte-operation sequence to
$\mathsf{RenderOriginal}$.  Dereferencing gives the same inputs to every
sealing, digest, MAC, validation, and output operation.  Coupled coins and
common oracles therefore give the same bytes, statuses, and state changes.
This step retains a write that became durable before returning an error;
the caller's status does not overwrite the persistence state.

An exact copy or pending-target application copies the object reference and
its original cache entry.  Aggregation collects the corresponding original
references, and resealing evaluates the new recipe using those references.
Each operation preserves the session-record handles and their save times.
It follows by induction on the recipe constructors that
$\mathsf{paths}(h)$ is the same in both machines.  A public hash or MAC
constructor contributes its common output and has an empty record-acquisition
path set.

On successful capture, output the cached object bytes in both experiments
and append its immutable descriptor to the adversary archive.  For each
version in its support, emit the same acquisition event at the actual
capture time.  Preserve the ordinary error for an unavailable object.
When a matching wrapping key is disclosed, use the common key identity and
reevaluate the enabled paths in the complete adversary archive.  For an
opening response, obtain $X_v$ from $\mathsf{RevealOriginal}$ on its first
use and from the wrapper's cache thereafter.  Every such call has
$v\in\mathcal L_T$.  Shared-key aliases and later key acquisition are thus
handled without another copy of $X_v$.  The optional key-extraction extension
preserves the invariant by induction on the least-fixed-point rounds.

Host deletion removes only host references.  The adversary archive and
already returned payloads persist.  Crash and resume use the same retained
object identifiers; neither operation recreates a reservation or changes a
save event.  Ordinary protocol actions are performed by the same
$\mathcal P$ with the same results.  This completes the induction, including
adaptive choices based on any preceding response.

For quantum private state, identical classical responses and the same public
oracle transformations induce the same transformation of that state at
each interaction.  The argument introduces no measurement of a quantum
oracle query and no oracle-dependent advice.  Finally, caching each
plaintext opening by $v$ proves the query bound; reservation and fresh-seal
indices were preserved separately throughout.
\end{proof}

The lemma places multiplicity in the correct ledger.  An adversary may retain
many ciphertext aliases of one version, and those ciphertexts may have many
different public heads.  All such byte-level observations remain in the view.
The repeated payload uses one original-record opening, whereas each genuinely
fresh seal and each public hash or MAC computation keeps its own resource
accounting.

% Manuscript fragment. Requires amsmath, amssymb, amsthm.
\subsection{Explicit guarded normalization of record readability}
\label{sec:guarded-provenance}

The path and readability definitions above already account for copying and
later key disclosure. We now give their explicit guarded normalization and
an executable realization for a record-decryption specialization. Here key
derivation has fixed public edges independent of recovered plaintext. The
optional plaintext-dependent key-extraction extension remains as defined
above and is outside this specialization. The set $\Omega_{\rm src}$ below
contains reservation-source identities; it is distinct from saved record
versions and from the original-payload table $\mathcal O[v]$. Likewise,
$\mathcal E^{\rm src}_t$ projects derivability onto protected source values,
whereas $\mathcal C_t$ and $\mathcal L_t$ track acquired and readable versions.
The graph is a fixed representation of admitted retained objects, not a
replacement for the provenance transition rules or the byte renderer.


A source identifier names one original sampling operation.  An acquisition
history records what the adversary obtains, including objects obtained at
different times.  These are distinct objects of the model: a copy preserves
the former, while an additional acquisition can change the latter.

\paragraph{Immutable record graph.}
Let $\mathcal R$ be a finite set of immutable encrypted record objects.
Each $r\in\mathcal R$ has a matching decryption-key identity $\kappa(r)$,
a set of embedded child records $\operatorname{child}(r)$, and a set
$\operatorname{orig}(r)$ of protected source values in its plaintext.
Containment is acyclic.  Record identities distinguish ciphertext objects;
source identities distinguish original sampling operations.  Several
records can contain the same source.  A rewrapped record has a new record
identity and may have a new key, while retaining its source identity.
Key identities are equivalence classes of equal actual keys. Equal key bytes
therefore have one identity across separate API objects.

Let $\mathcal D\subseteq\mathcal K\times\mathcal K$ be fixed public key
derivation edges.  An edge $(k,k')$ means that possession of $k$ permits
deterministic computation of $k'$.  This interface represents, for example,
derivation of an operation-wrapping key from the journal key.  In the
theorem below, these edges do not depend on plaintext recovered from records.
Write $D^*$ for their reflexive transitive closure.

Let $\Omega_{\rm src}=\bigcup_{r\in\mathcal R}\operatorname{orig}(r)$ be the
declared source domain.  At a trace prefix $t$, let $A_t\subseteq\mathcal R$ be all record objects
ever acquired by that prefix, $K_t\subseteq\mathcal K$ all directly
disclosed keys, and $\mathcal S^{\rm dir}_t\subseteq\Omega_{\rm src}$ all directly disclosed source values, represented
by their source identities.  These are cumulative sets.  The decryption-key
closure is
\[
  K_t^*=\{k':\exists k\in K_t,\ (k,k')\in D^*\}.
\]
The concrete derivation rules are
\[
 \frac{k\quad(k,k')\in\mathcal D}{k'},\qquad
 \frac{r\quad\kappa(r)}{c}\ (c\in\operatorname{child}(r)),\qquad
 \frac{r\quad\kappa(r)}{o}\ (o\in\operatorname{orig}(r)).
\]
Their least closure starts from $A_t,K_t,\mathcal S^{\rm dir}_t$.  A fact $o$ in the source
sort means that the protected value of source $o$ is derivable; it does
not mean that a public identifier has become known.  The source projection
of this closure is denoted $\mathcal E^{\rm src}_t$.

\paragraph{Guarded routes.}
For every path $p=(r_0,\ldots,r_\ell)$ in the containment graph ending
at a record with $o\in\operatorname{orig}(r_\ell)$, define
\[
 g(p)=\{\kappa(r_0),\ldots,\kappa(r_\ell)\}.
\]
The normalized table contains a route $(r_0,g(p),o,p)$ for every such
path.  Its exposure expression is
\[
 \widehat{\mathcal E}^{\rm src}_t=\mathcal S^{\rm dir}_t\ \cup\
 \{o:\exists(r,g,o,p),\ r\in A_t\ \land\ g\subseteq K_t^*\}.
 \tag{GP}
\]
Alternative records and paths are disjoined, while all keys on one path
are conjoined.  Equal key identities occur only once in the guard set.

\begin{theorem}[Exact guarded normalization]
\label{thm:guarded-normalization}
For every finite record graph and derivation interface above, and every
acquisition prefix $(A_t,K_t,\mathcal S^{\rm dir}_t)$, the concrete and normalized exposure
sets agree: $\mathcal E^{\rm src}_t=\widehat{\mathcal E}^{\rm src}_t$.  The equality applies to all prefixes of
an adaptively chosen history.
\end{theorem}
\begin{proof}
Key derivations depend only on disclosed keys and the fixed relation
$\mathcal D$, so their least closure is exactly $K_t^*$.
For the forward inclusion, take a source $o\in \mathcal E^{\rm src}_t$.
If it was directly disclosed then $o\in \mathcal S^{\rm dir}_t$.
Otherwise its first derivation decrypts a record $r_\ell$ with the
matching key.  If that record was not acquired directly, its first
availability follows by decrypting a parent.  Follow these first
derivations backwards.  The finite acyclic record graph makes this
process terminate at an acquired record $r_0$.  Every intervening key
lies in $K_t^*$, yielding a route satisfying (GP).
For the reverse inclusion, take a satisfied route.  Starting with its
acquired first record, successively decrypt along the path.  Every
required key is in $K_t^*$, so each child and finally $o$ is in the
concrete least closure.  Direct disclosures are in that closure by
definition.  Both inclusions hold for each fixed prefix; no independence
or advance choice of the acquisition schedule was used.
\end{proof}

\paragraph{Efficient evaluation.}
Equation (GP) is a declarative normal form: a compressed record DAG can
have exponentially many paths.  An efficient leakage evaluator instead
computes the iterative decryption and key-derivation fixed point, without
enumerating those paths.  Let $N$ be the total number of record, key and
source atoms, and let $M$ count rule and payload incidences.  There are at
most $N$ strict-growth rounds; repeated rule scans therefore take
$O(N(N+M))$ time and $O(N+M)$ space under standard indexed set and
dictionary representations, with source-dictionary operations included
in this accounting.  The finite checker materializes the route table only
for its declared small fixtures, to compare that normal form against the
forward evaluator.

\begin{corollary}[Copies, deletion and late disclosure]
\label{cor:guarded-history}
Exposure is monotone under additional acquisitions, direct disclosures
and key disclosures.  Copying or rewrapping a source adds routes to the
same source identity.  Local deletion of a record previously acquired
does not remove its routes from adversarial knowledge.  In particular,
a later matching-key disclosure can expose a source in a ciphertext
acquired before that disclosure and deleted locally in the meantime.
\end{corollary}
\begin{proof}
The cumulative sets and their key closure only grow, and every predicate
in (GP) is positive.  A source-preserving copy changes a route's record
identity or key guards but not its terminal source.  Local deletion
restricts future acquisition from that local store; it removes no member
of $A_t$ and therefore cannot invalidate a satisfied route.
\end{proof}

\paragraph{Interface significance.}
The theorem gives an exact finite normalization of the chosen
record-decryption semantics.  It supports a precise original-operation
leakage oracle: a query for source $o$ can be evaluated against the
complete acquisition history, with a path and its key guards as a witness.
The cryptographic transition from native ciphertexts and public
commitments to this interface has a separate simulation obligation.
The derivation rules do not model guessing, cryptanalysis, timing leakage
or a quantum decryption oracle.  Accordingly, absence from $\mathcal E^{\rm src}_t$ is
recorded as non-derivability by these rules; message indistinguishability
outside an exposure set is supplied by a security theorem for the
surrounding protocol.

\paragraph{Native shape and finite checks.}
At the fixed research source, operation reservations are encrypted under
a wrapping key derived from the journal key, and tokens are retained
within encrypted durable images or pending targets.  The corresponding
route therefore includes the image key and the token key, with their
actual derivation edge.  A fresh aggregate image can preserve an older
token and its original source.  Public image-head digests are distinct
observations from the record objects needed by the decryption rules.
The executable supplement independently evaluates iterative decryption
and guarded-route normalization, exhausts all acquisition/key/direct-source
subsets for its declared finite graphs, and checks histories with late
key disclosure and local deletion.  Its counts and source hash are in
the accompanying machine-readable result; those checks validate the
implemented finite model and supplement the general proof above.

% Standalone manuscript fragment; no preamble.
% Requires proposition and proof environments (e.g. amsthm).
\paragraph{Why exposure must retain conjunctive guards.}
The following elementary separation explains the conjunction in (GP).

\begin{proposition}[Acquisition exposure is not additive]
\label{prop:acquisition-nonadditivity}
Let $r$ be a retained-record atom, $k$ a wrapping-key atom, and $s$ a
protected-seed atom associated with an origin identifier $o$.
For $A\subseteq\{r,k\}$, let $\operatorname{Cl}(A)$ be the least set
containing $A$ and closed under the rule
\[
  \{r,k\}\;\Longrightarrow\;s.
\]
There are no other rules in this fragment.  Define the exposed-origin set by
\[
  E(A)=\{\,o : s\in\operatorname{Cl}(A)\,\}.
\]
There is no fixed function
$f:\{r,k\}\longrightarrow\mathcal{P}(\{o\})$ satisfying
\[
  E(A)=\bigcup_{x\in A} f(x)
  \qquad\text{for every }A\subseteq\{r,k\}.
\]
\end{proposition}

\begin{proof}
The closure rule gives
\[
  E(\{r\})=E(\{k\})=\varnothing,
  \qquad
  E(\{r,k\})=\{o\}.
\]
Exactness for the two singleton acquisitions forces
$f(r)=f(k)=\varnothing$.  Their union is then empty for the joint
acquisition, contradicting $E(\{r,k\})=\{o\}$.
\end{proof}

\paragraph{Meaning of exposure.}
Membership of an origin $o$ in $E(A)$ records derivability of its
protected seed under the specified closure rules; the identifier itself
may be public. Computational estimation is measured separately by success
probability. A one-time-pad illustration makes the dependence on joint
observations explicit: let $S,W$ be independent uniform $\lambda$-bit strings and
$C=S\mathbin{\oplus}W$, with no additional observation depending on $S$ or
$W$.  Then
\[
  H(S\mid C)=H(S\mid W)=\lambda,
  \qquad H(S\mid C,W)=0.
\]
From either observation alone, any estimator of $S$ succeeds with
probability at most $2^{-\lambda}$; from both, $S=C\mathbin{\oplus}W$
recovers it exactly. The XOR construction serves here as an
information-theoretic illustration of acquisition; the protocol's record
encryption is handled by its cryptographic record interface.

\paragraph{Consequence for provenance routes.}
A source label and an exposure decision serve different purposes.  Copying
a retained record preserves the source label, while subsequent acquisition
of a matching key may change exposure.  A guarded-route representation
retains this distinction: alternative acquired records supply alternative
routes to the same source, and every route retains the conjunction of its
required keys.  Exposure is therefore evaluated from the joint acquisition
and key-derivation history. The separation applies specifically to an
exact exposure representation formed by the union of fixed per-object
labels.


\subsection{A recovery experiment with acquisition-based eligibility}

\begin{definition}[Fixed-target recovery experiment]
\label{def:recovery-game}
Fix a protocol profile $\mathcal P$ and finite execution budgets, together with
a designated session $s$, recovery cut $c$, target epoch $e$, and candidate
generation and encapsulation reservation handles $g_e,a_e$.  Handles may be
named by original invocation slots before those invocations occur; a missing
invocation will fail the reservation-event test below.  Public
protocol messages may be delayed, duplicated, reordered, or replaced.
The adversary can capture the declared retained objects, reveal wrapping
keys, reveal actual signing keys or signer files, and reveal designated
reservation tapes.  The cut exposes the selected wrapping keys; preceding
roots may already be public or disclosed by the stated prior captures.
These disclosures have their exact scope in the event ledger.

For one fresh logical application send, the adversary supplies equal-length
messages $(m_0,m_1)$.  The challenger samples $b\leftarrow\{0,1\}$ and runs
the ordinary sending transition on $m_b$.  Rejection or an unknown outcome
keeps its ordinary protocol behavior.  A retained accepted send is identified
by $(s,e,d,j)$, including direction $d$ and the epoch-scoped logical index $j$;
exact retries reuse its original ciphertext.  The challenge is valid only
when that logical send is accepted in the designated epoch.  Subsequent
protocol actions and disclosures continue until the adversary outputs $b'$.
\end{definition}

Write $\iota_A,\iota_B$ for the two actual signing-key identities configured
for the target session.  The predicate
$\mathsf{ReservedFor}(s,e,g_e,a_e)$ checks the occurrence of the named honest
generation and encapsulation reservation events in the target epoch's local
durable plans, with their target session and roles.  It checks event metadata,
without asserting that their outputs form the accepted target exchange.
Eligibility uses the following event conditions over the complete execution:
\begin{align}
 \mathsf{Eligible}_{s,e,c}(\tau)\ \Longleftrightarrow\quad&
 \mathsf{ValidTarget}(\tau)
 \ \wedge\ \mathsf{ReservedFor}(s,e,g_e,a_e)
 \nonumber\\[-1mm]
 &{}\wedge\ c<\mathsf{res}(g_e)
          \ \wedge\ c<\mathsf{res}(a_e)
 \nonumber\\
 &{}\wedge\ \iota_A,\iota_B
                  \notin\mathsf{IdentityExposed}(\tau)
 \nonumber\\
 &{}\wedge\ \mathsf{EntropyExposed}(s,e;\tau)=\varnothing
 \nonumber\\
 &{}\wedge\ \nexists v\in\mathcal C_T:
       \mathsf{subject}(v)=s\ \wedge\ c<\mathsf{save}(v).
 \label{eq:eligibility}
\end{align}
Here $\mathsf{EntropyExposed}(s,e;\tau)$ records the target epoch's designated
entropy disclosures.  The separate event $\mathsf{Attr}$ states that these
named reservation events form the accepted target's unique usable
generation/encapsulation pair, with an adapter-computable association to the
original slots.  In particular, it connects the generation public key, the
encapsulation input, the authenticated response, and the target's subsequent
key schedule.  This matching is a
source/authentication obligation: a failure of $\mathsf{Attr}$ can occur while
$\mathsf{Eligible}_{s,e,c}$ holds and is charged in the computational
accounting.

The final line of~\eqref{eq:eligibility} concerns every saved version of session
$s$, across its roles and epochs.  It uses acquisition
$\mathcal C_T$, including protected bytes, at the time they are obtained.
Readability $\mathcal L_T$ is evaluated separately. Other epochs' entropy disclosures retain their ordinary
behavior.  Both peers' confirmation is not an additional send prerequisite.

Define the unconditional centered advantage by
\begin{equation}
 \mathsf{Adv}^{\rm rec}_{\mathcal P}(\mathcal A)
   =
 \left|
   \mathbb E\!\left[
    \mathbf 1_{\mathsf{Eligible}_{s,e,c}(\tau)}
                       (-1)^{b\oplus b'}\right]
 \right|.
 \label{eq:recovery-score}
\end{equation}
Invalid or ineligible executions contribute zero.  This definition retains
failures and adaptive disclosures in the probability space and avoids
conditioning a reduction on an uncontrolled eligibility probability.
If every execution with $\mathsf{Eligible}_{s,e,c}=0$ is completed by an
independent fair guess, the conventional advantage over one half is
\[
 \mathsf{Adv}_{\rm msg}
   =\left|\Pr[\text{completed guess is correct}]-\tfrac12\right|
   =\tfrac12\,\mathsf{Adv}^{\rm rec}_{\mathcal P}(\mathcal A).
\]
The computational accounting uses this latter convention.

The adaptive version retains one application challenge and permits the
candidate selector to choose its coordinates and reservation-event handles
adaptively from the observed interaction. Eligibility
requires those coordinates to identify the challenge's accepted logical send;
selection never transfers its test bit to another send.  Eligibility
and the usable-pair attribution are evaluated on that experiment's complete
history.  An origin-pair reduction fixes a slot pair in $[G]\times[C]$ and
scores the event that the final selector attributes the target to that pair.
The selector and disclosure predicates are recomputed in each comparison
world.  No independence between that selection and the transcript is imposed.

The record restriction in~\eqref{eq:eligibility} is the one used by the
retained multi-epoch symbolic theorem: its \textsf{Snapshot} action outputs
a protected record, while \textsf{Cut} separately releases the wrapping keys.
The theorem establishes symbolic plaintext secrecy with both original
reservations after the cut, key-scoped identity exclusions, the target
entropy exclusion, and the stated post-cut saved-record acquisition
exclusion~\cite{qperiaptSymbolic2026}.  Its joint replay verifies the main
lemma and all 24 reusable dependencies.  Equation~\eqref{eq:recovery-score}
specifies the computational message experiment to which the cryptographic
modules are to be connected; the alias lemma establishes its record-interface
normalization.

\begin{corollary}[Safe stopping at a forbidden acquisition]
\label{cor:target-acquisition-stop}
Consider the fixed-target experiment of
Definition~\ref{def:recovery-game}.  Suppose an actual successful capture
contains a target-session version $v$ with $c<\mathsf{save}(v)$, identified
by the original-record metadata.  For any terminal bit $D$, replacing the
remainder of that execution by output zero leaves
\[
   \Pr[\mathsf{Eligible}_{s,e,c}(\tau)\wedge D=1]
\]
unchanged.  Replacing its terminal score by zero also leaves the expectation
inside~\eqref{eq:recovery-score} unchanged.

In particular, if a captured target-session version retains the selected
reservation $o\in\{g_e,a_e\}$, then
$c<\mathsf{res}(o)$ and~\eqref{eq:origin-causality} establish this stopping
condition.  The wrapping key may be obtained before or after the capture.
\end{corollary}

\begin{proof}
The successful capture inserts $v$ into $\mathcal C_t$ immediately.
Its save time and subject are immutable, and acquisition remains recorded
after host deletion, resealing, confirmation, or process loss.  Thus the
last conjunct of~\eqref{eq:eligibility} is false for every continuation.
Each such continuation already contributes zero to both quantities.
On executions that never reach this condition, the experiment is unchanged.
The final assertion follows from
$c<\mathsf{res}(o)\leq\mathsf{save}(v)$.
\end{proof}

The stopping test uses the successful acquisition event and its record
metadata; it can run before the response's private bytes are serialized by
a challenge adapter.  It is not an SDK rejection of the capture.  A public
digest, an unrelated record, or an unsuccessful capture keeps its real
behavior.  In particular, difficulty computing a permitted digest or report
MAC does not trigger this corollary.
For adaptive target selection, stopping is applied to the corresponding
fixed-candidate event in the reduction.  One possible candidate becoming
ineligible does not justify terminating a game that can still select a
different eligible target.

Two short executions illustrate why both kinds of origin are necessary.
If $v$ was saved before $c$ and a sibling write reseals it after $c$, the
acquisition still refers to the original pre-cut $v$.  The new global image
time alone does not falsify the final conjunct of~\eqref{eq:eligibility}.
If instead $v$ was saved after $c$, a later write attributed to a sibling
does not change its subject or make its capture eligible.  Likewise, a
pre-cut generation reservation retained through a delayed receipt remains
pre-cut when its KEM computation completes later. The reservation condition
in~\eqref{eq:eligibility} determines that case, with the endpoints retaining
their respective confirmed counters.

\subsection{Connection to the cryptographic argument}

The normalization theorem fixes the original objects whose bytes, acquisition,
and key identities the reduction must represent. The next section gives
constructive image/intent and report interfaces for a selected KEM operation
whose public outputs are available while its original random tape is opaque.
Their algorithms preserve original-object caches, public commitments, exact
retries, and actual durable outcomes. The fixed-target stopping rule above
extends to an adaptive provenance certificate in
Lemma~\ref{lem:adaptive-origin-stop}.

Usable attribution is computed from available semantic metadata. Original
generation and encapsulation slots belong to a fixed lifetime budget
$[G]\times[C]$: a fresh draw receives a new slot, while copies retain their
source. Each candidate induces its own protected records and report-key
closure. Proposition~\ref{prop:native-local-accounting} decomposes the score
in the native experiment before comparing those local continuations. It
retains both unattributable eligible executions and permitted key projections
as explicit native-game terms.

The computational composition evaluates accepted operations, eligibility, and
all public outputs anew in each comparison world. Exact component repeats
retain one response and one correctness event. The native constructor map in
Section~\ref{sec:native-correspondence} identifies the storage operations used
by these interfaces; Section~\ref{sec:discussion} specifies the joint root and
message continuation to which their bounds connect.
